\section{Supplementary statements, dependencies, and open problems}
\label{sec:supplementary-theorems}

This section is a \emph{cross-reference ledger} for statements that are mentioned elsewhere in the manuscript.
To avoid circularity and accidental promotion of heuristic physics to mathematics, the items below are labeled by status:
\begin{itemize}
	\item \textbf{[PROVED]}: proved in this manuscript (with an explicit proof).
	\item \textbf{[LITERATURE]}: a standard theorem with a precise reference.
	\item \textbf{[HYPOTHESIS]/[OPEN]}: used only conditionally or posed as an open problem.
\end{itemize}
Where relevant, we also point to Appendix~\ref{sec:open-problems} for a consolidated list.

\subsection{Items anchored in strong coupling}

\begin{theorem}[Strong-coupling analyticity of the free energy \textnormal{[PROVED]}]
\label{thm:convex-analytic}
There exists $\beta_0>0$ such that the free energy density $f(\beta)$ of $SU(N)$ lattice Yang--Mills theory is analytic for $|\beta|<\beta_0$.
\end{theorem}

\begin{theorem}[Strong-coupling string tension lower bound \textnormal{[PROVED]}]
\label{thm:sigma-positivity-final}
There exists $\beta_0>0$ such that for $0<\beta<\beta_0$ the string tension $\sigma(\beta)$ is strictly positive.
\end{theorem}

\subsection{Statements that remain conditional}

\begin{theorem}[Global analyticity of the free energy \textnormal{[OPEN]}]
\label{thm:global_analyticity_supp}
\emph{Open problem.} Prove that $f(\beta)$ is real-analytic for all $\beta>0$ (or, more weakly, that there is no bulk transition obstructing the strong-coupling regime from being continuously connected to the weak-coupling regime).
\end{theorem}

\begin{theorem}[Non-Abelian correlation inequalities / monotonicity \textnormal{[OPEN]}]
\label{thm:nonabelian_monotonicity_supp}
\emph{Open problem.} Establish a rigorous monotonicity principle for Wilson loop expectations or string tension in non-Abelian $SU(N)$ lattice gauge theory at all couplings.
In particular, no argument in this manuscript relies on a GKS/Griffiths-type inequality for $SU(N)$.
\end{theorem}

\begin{theorem}[Uniform log-Sobolev inequality \textnormal{[OPEN]}]
\label{thm:uniform-lsi-induction}
\emph{Open problem.} Prove a uniform (in volume and along the scaling trajectory) log-Sobolev inequality for the lattice gauge measure, or provide an alternative route to uniform mixing/correlation-length control.
\end{theorem}

\begin{theorem}[Cluster decomposition / finite correlation length at all couplings \textnormal{[OPEN]}]
\label{thm:cluster}
\emph{Open problem.} Prove finite correlation length (or exponential clustering) for all $\beta$ in infinite volume.
\end{theorem}

\begin{theorem}[Dimensionless ratio bound $R(\beta)$ \textnormal{[OPEN]}]
\label{thm:ratio-bound}
\emph{Open problem.} Prove a uniform positive lower bound for $R(\beta)=\Delta(\beta)/\sqrt{\sigma(\beta)}$ along a scaling trajectory.
\end{theorem}

\begin{theorem}[Giles--Teper-type inequality \textnormal{[OPEN]}]
\label{thm:giles-teper_supp}
\emph{Open problem.} Establish a rigorous inequality of the form $\Delta\ge c_N\sqrt{\sigma}$ with explicit $c_N>0$ (or any nontrivial quantitative relation between gap and string tension) from first principles.
\end{theorem}

\begin{theorem}[Stochastic quantization SPDE well-posedness \textnormal{[OPEN]}]
\label{thm:spde-wellposed}
\emph{Open problem.} Prove global well-posedness and identification of an invariant measure for a mathematically controlled stochastic quantization of 4D Yang--Mills.
\end{theorem}

\begin{theorem}[Osterwalder--Schrader axioms in the continuum limit \textnormal{[OPEN]}]
\label{thm:full-os}
\emph{Open problem.} Construct a continuum scaling limit satisfying OS axioms (and then reconstruct a Wightman QFT) from the lattice measures.
\end{theorem}

\section{Factorization (programmatic)}
\label{sec:factorization}
This section records definitions/targets used in programmatic discussions.

\begin{definition}[Epsilon Factorization]
\label{def:eps-fact}
The measure factorizes up to an error $\epsilon$ at scale $L$.
\end{definition}

\begin{theorem}[Wilson loop convergence \textnormal{[OPEN]}]
\label{thm:wilson-convergence}
\emph{Open problem.} Prove convergence of Wilson loop expectations along a continuum scaling trajectory.
\end{theorem}

\begin{theorem}[Wightman axioms \textnormal{[OPEN]}]
\label{thm:wightman}
\emph{Open problem.} After OS reconstruction, verify Wightman axioms for the reconstructed theory.
\end{theorem}

\begin{theorem}[Equicontinuity / tightness targets \textnormal{[OPEN]}]
\label{thm:equicontinuity}
\emph{Open problem.} Prove the uniform equicontinuity/tightness estimates needed for compactness arguments.
\end{theorem}

\begin{theorem}[Pure spectral gap \textnormal{[OPEN]}]
\label{thm:pure-spectral-gap}
\emph{Open problem.} Prove a nonzero spectral gap in the continuum limit for 4D pure Yang--Mills.
\end{theorem}

\begin{theorem}[String tension positivity in the continuum \textnormal{[OPEN]}]
\label{thm:sigma-positive-continuum}
\emph{Open problem.} Prove nonzero string tension in the continuum scaling limit.
\end{theorem}

% (OS axioms are recorded above as an open problem.)

\section{Gap closure (programmatic cross-reference)}
\label{app:definitive-gap-closure}
This label is retained for cross-referencing legacy text.
In the rigorous-scope version of the manuscript, any content labeled as a \emph{definitive gap closure} is treated as programmatic unless fully proved.

\begin{corollary}[Geometric Positivity]
\label{cor:h-geom-positive}
The geometric Hamiltonian is strictly positive.
\end{corollary}

\section{Wilsonian QCD}
\label{sec:wilsonian-qcd}
This section discusses the Wilsonian approach to QCD.

\section{String tension (programmatic)}
\label{sec:string-tension}
This section analyzes the string tension.

\begin{theorem}[Complete rigorous proof \textnormal{[NOT CLAIMED]}]
\label{thm:complete-rigorous}
This label is retained for backward compatibility with earlier drafts.
No complete rigorous proof claim is made in this version.
\end{theorem}

\begin{theorem}[Large Field Bounds]
\label{thm:large-field}
The large field configurations are suppressed by the action.
\end{theorem}

\begin{theorem}[Main theorem \textnormal{[NOT CLAIMED]}]
\label{thm:main}
This label is retained for backward compatibility with earlier drafts.
An unconditional mass gap theorem is not claimed in this version.
\end{theorem}

\section{Variance Transport}
\label{sec:variance-transport-sec12}
This section discusses variance transport.

\begin{theorem}[Basic Transfer Matrix]
\label{thm:transfer-basic}
The transfer matrix is well-defined and bounded.
\end{theorem}

\section{Introduction}
\label{sec:intro}
Introduction to the problem.

\section{Spectral Stability}
\label{sec:spectral-stability}
This section discusses the stability of the spectrum.

\begin{theorem}[Universality]
\label{thm:universality}
The continuum limit is universal and independent of the lattice regularization.
\end{theorem}

\section{Universality}
\label{sec:universality}
This section discusses the universality of the continuum limit.

\section{Axiomatic Mass Gap}
\label{sec:axiomatic-mass-gap}
This section discusses the mass gap from an axiomatic perspective.

\begin{theorem}[Spectral Wilson Loop]
\label{thm:spectral-wilson}
The spectral properties of the Wilson loop operator are related to the mass gap.
\end{theorem}

\begin{theorem}[Dimensional Transmutation Tools]
\label{thm:dimensional-transmutation-tools}
The tools for analyzing dimensional transmutation include the beta function and the renormalization group equation.
\end{theorem}


